\documentclass[11pt]{article}
\usepackage{series}
\usepackage{amsmath,amssymb,amsthm}
\usepackage{geometry}
\usepackage{hyperref}
\usepackage{enumitem}

\geometry{margin=1in}


% ---------- metadata ----------
\title{The Modal Triplet Theory Program D1:\\ The Dark Sector as Missing Encodings
in the Modal Triplet Theory Program}
\author{Peter Nero}
\date{January, 2026}

\begin{document}

\maketitle

\begin{abstract}
We reinterpret dark matter and dark energy as phenomenological regimes arising
from restricted availability of admissible encodings in the Modal Triplet Theory
(MTT) framework. No new obstruction types or encoding responses are introduced.
Instead, dark sector phenomena are shown to emerge when local descriptive
encodings fail or decouple, leaving only relational or statistical encodings
active.

Dark matter is associated with regimes dominated by statistical ensemble
encodings (E9), in which matter influences kinematic consistency and curvature
through aggregate effects but lacks stable local descriptive structure. Dark
energy is associated with regimes dominated by relational encodings (E8), in
which global transition structure persists while local matter descriptions are
absent or ineffective.

This perspective explains the universality, gravitational coupling, and apparent
non-interaction of the dark sector without postulating new particles, forces, or
fields. The dark sector is not fundamental; it reflects descriptive incompleteness
rather than new ontology. Observational consequences are discussed at the level
of effective behavior, without modifying the structural or encoding-level results
of the theory.
\end{abstract}

\tableofcontents
\newpage

% ============================================================
\section{Introduction and Scope}

Observations in cosmology and astrophysics indicate the presence of gravitational
effects not accounted for by visible matter and radiation. These effects are
commonly attributed to dark matter and dark energy, often modeled as new
particle species, exotic fluids, or modifications of gravitational dynamics.
Despite extensive effort, no direct non-gravitational detection of dark sector
degrees of freedom has been confirmed.

The purpose of the present paper is to reinterpret these phenomena within the
Modal Triplet Theory Program. Rather than introducing new encodings, obstruction
types, or dynamical laws, we show that dark sector behavior arises naturally as a
consequence of \emph{missing or ineffective encodings} in certain regimes of
admissible description.

In the MTT framework, reduced descriptions are local and encoding-dependent.
Multiple descriptive encodings (E1--E9) may be admissible in different regimes,
and no single encoding is guaranteed to be globally available. When local
descriptive encodings (such as operator, spectral, or field-based charts) fail
or decouple, only relational or statistical encodings may remain admissible.
These regimes manifest observationally as dark sector phenomena.

Specifically, we argue that:
\begin{itemize}
\item \emph{Dark matter} corresponds to regimes dominated by statistical ensemble
encoding (E9), in which aggregate distributions influence kinematic consistency
and curvature without stable local degrees of freedom.
\item \emph{Dark energy} corresponds to regimes dominated by relational encoding
(E8), in which global transition structure persists while local matter
descriptions are absent or dynamically ineffective.
\end{itemize}

In this interpretation, the dark sector is not a new form of matter or energy,
but a manifestation of descriptive incompleteness. It couples gravitationally
because gravity encodes kinematic consistency universally, but it lacks gauge
and local interaction signatures because those encodings are not active.

\begin{remark}[Interpretive scope]
This paper introduces no new encoding responses and no new structural
assumptions. It interprets existing observations as consequences of restricted
encoding availability. The structural and encoding-level results of the Modal
Triplet Theory Program remain unchanged.
\end{remark}

The organization of the paper is as follows. In Section~2 we summarize the
encoding landscape and clarify what it means for an encoding to be ``missing''
or ineffective. Section~3 analyzes dark matter as a statistical-encoding-dominated
regime. Section~4 analyzes dark energy as a relational-encoding-dominated regime.
Section~5 discusses observational consequences and constraints. Section~6
addresses common objections and alternative interpretations. Section~7
summarizes the results and places the dark sector interpretation within the
broader MTT program.

\begin{remark}[Position in the program]
This paper belongs to the D-layer of the Modal Triplet Theory Program. It is
interpretive and phenomenological. Readers interested only in structural
necessity or encoding responses may omit this paper without loss of logical
completeness.
\end{remark}

% ============================================================
% Next section will be:
% \section{Encoding Availability and Descriptive Regimes}
% ============================================================
\section{Encoding Availability and Descriptive Regimes}

In this section we formalize the notion of encoding availability and explain how
different descriptive regimes arise within the Modal Triplet Theory framework.
The dark sector interpretation depends crucially on the fact that not all
admissible encodings are available or effective in all regimes.

\subsection{Encoding availability}

The Modal Triplet Theory Program allows multiple descriptive encodings
(E1--E9) to coexist in principle, but does not require them to be simultaneously
admissible everywhere.

\begin{definition}[Encoding availability]
An encoding is said to be \emph{available} in a region if it admits an admissible
reduced description satisfying locality, overlap consistency, and refinement
stability in that region.
\end{definition}

Encoding availability is therefore a regime-dependent property, not a global
feature of the theory.

\subsection{Encoding failure and decoupling}

Encodings may become unavailable in two distinct ways.

\begin{definition}[Encoding failure]
An encoding fails in a region if no admissible reduced description of that type
exists there.
\end{definition}

\begin{definition}[Encoding decoupling]
An encoding decouples in a region if it remains formally admissible but does not
contribute dynamically or observationally to kinematic persistence or
consistency.
\end{definition}

Both mechanisms lead to effective absence of an encoding in that regime.

\subsection{Hierarchy of descriptive encodings}

Different encodings carry different descriptive power.

\begin{remark}
Local operator, field-theoretic, or spectral encodings (E1--E6) require strong
regularity and stability conditions. Relational (E8) and statistical (E9)
encodings are weaker: they persist even when local structure is degraded or
unavailable.
\end{remark}

This hierarchy explains why E8 and E9 often dominate in extreme or poorly
resolved regimes.

\subsection{Descriptive regimes}

We now define descriptive regimes.

\begin{definition}[Descriptive regime]
A \emph{descriptive regime} is a region of admissible description characterized
by the subset of encodings that remain available or effective.
\end{definition}

Different regimes may admit different subsets of E1--E9.

\subsection{Transition between regimes}

Transitions between descriptive regimes are generic.

\begin{remark}
As admissibility conditions change—for example due to scale, refinement,
instability, or horizon effects—encodings may become unavailable or decouple.
Such transitions do not require new physics; they reflect changes in descriptive
capacity.
\end{remark}

These transitions are central to the interpretation of the dark sector.

\subsection{Gravity as universal bookkeeping}

Gravity remains active across all regimes.

\begin{remark}
Gravity encoding enforces kinematic consistency universally. It remains operative
even when other encodings fail or decouple. This universality explains why dark
sector phenomena couple gravitationally in all observed regimes.
\end{remark}

By contrast, gauge and local matter encodings may become ineffective.

\subsection{Dark sector as a regime, not an entity}

We emphasize the interpretive stance.

\begin{remark}
The dark sector is not identified with a new substance, field, or interaction.
It is identified with descriptive regimes in which only relational (E8) or
statistical (E9) encodings remain available, while local descriptive encodings
are absent or dynamically irrelevant.
\end{remark}

This reframing eliminates the need for additional ontology.

\subsection{Preview: dark matter and dark energy}

We now prepare the specific analyses.

\begin{remark}
In the next section we analyze dark matter as a regime dominated by statistical
ensemble encoding (E9), while in Section~4 we analyze dark energy as a regime
dominated by relational encoding (E8).
\end{remark}

% ============================================================
% Next section will be:
% \section{Dark Matter as Statistical-Encoding Dominance}
% ============================================================
\section{Dark Matter as Statistical-Encoding Dominance}

In this section we interpret dark matter phenomena as arising from descriptive
regimes dominated by statistical ensemble encoding (E9). In such regimes,
aggregate structure remains admissible and gravitationally active, while local
descriptive encodings are absent, unstable, or dynamically ineffective.

\subsection{Statistical ensemble encoding (E9)}

We briefly recall the role of E9 encodings.

\begin{remark}
Statistical ensemble encoding (E9) describes systems in terms of aggregate or
distributional properties rather than local degrees of freedom. It remains
admissible when detailed local structure cannot be stably resolved, provided
sufficient statistical regularity persists.
\end{remark}

E9 does not require:
\begin{itemize}
\item localized particles;
\item well-defined local fields;
\item operator or spectral descriptions.
\end{itemize}

It requires only stable ensemble-level behavior.

\subsection{Loss of local descriptive encodings}

Dark matter regimes are characterized by loss or decoupling of local encodings.

\begin{remark}
In many astrophysical and cosmological regimes, local operator, field-theoretic,
or particle-based encodings (E1--E6) may fail due to instability, coarse-graining,
or lack of admissible refinement. When this occurs, only statistical encodings
remain available.
\end{remark}

The absence of local encodings does not imply absence of structure, but absence
of local describability.

\subsection{Gravitational activity of statistical encodings}

Statistical encodings remain gravitationally active.

\begin{lemma}
In a regime dominated by statistical encoding (E9), aggregate distributions
contribute to kinematic consistency encoding and therefore source gravitational
effects.
\end{lemma}

\begin{proof}
Gravity encoding enforces kinematic consistency universally, independent of the
descriptive encoding used. Statistical encodings contribute aggregate structure
that affects admissible continuation and curvature bookkeeping, and therefore
enter gravitational consistency conditions.
\end{proof}

This explains why dark matter influences gravitational dynamics despite lacking
local interactions.

\subsection{Absence of non-gravitational interactions}

Statistical dominance explains the non-interacting character of dark matter.

\begin{remark}
Gauge interactions require local descriptive encodings with well-defined
redundancy structure. In regimes where only statistical encoding is available,
such redundancy bookkeeping cannot be implemented locally. As a result,
non-gravitational interactions are suppressed or absent.
\end{remark}

This explains the failure of direct detection experiments without invoking new
forces or symmetries.

\subsection{Effective mass and clustering}

Dark matter exhibits effective mass and clustering behavior.

\begin{remark}
Statistical encodings admit stable aggregate distributions that can cluster under
gravitational consistency encoding. These distributions behave effectively as
mass sources, despite lacking particle interpretation.
\end{remark}

Thus dark matter appears massive and clumped without being particulate.

\subsection{Scale dependence}

Statistical dominance is scale-dependent.

\begin{remark}
At sufficiently small scales, local encodings may re-emerge and statistical
encoding may break down. At sufficiently large scales, statistical encoding
becomes dominant. This explains why dark matter effects are prominent at galactic
and cosmological scales but suppressed at laboratory scales.
\end{remark}

No new scale-dependent physics is required.

\subsection{Comparison with particle dark matter}

We contrast this interpretation with particle models.

\begin{remark}
Particle dark matter models introduce new local degrees of freedom to explain
gravitational anomalies. The statistical-encoding interpretation explains the
same phenomena without introducing new ontology, attributing them instead to the
absence of local descriptive encodings.
\end{remark}

The two approaches make different predictions regarding non-gravitational
interactions.

\subsection{Structural advantages}

The statistical-encoding interpretation has several advantages.

\begin{itemize}
\item It explains universal gravitational coupling naturally.
\item It explains non-detection in non-gravitational channels.
\item It avoids fine-tuning of particle properties.
\item It aligns with admissibility and refinement constraints.
\end{itemize}

These features follow structurally, not by design.

\subsection{Interpretive caution}

We emphasize the scope of the claim.

\begin{remark}
This interpretation does not exclude the existence of new particles or fields in
nature. It asserts only that dark matter phenomena do not require them at the
structural level. Statistical-encoding dominance provides a sufficient and
economical explanation within the Modal Triplet Theory framework.
\end{remark}

\subsection{Preview: dark energy}

Dark matter and dark energy arise from different descriptive regimes.

\begin{remark}
In the next section we analyze dark energy as a regime dominated by relational
encoding (E8), where global transition structure persists while even statistical
local structure is ineffective.
\end{remark}

% ============================================================
% Next section will be:
% \section{Dark Energy as Relational-Encoding Dominance}
% ============================================================
\section{Dark Energy as Relational-Encoding Dominance}

In this section we interpret dark energy phenomena as arising from descriptive
regimes dominated by relational encoding (E8). In such regimes, global transition
structure remains admissible while local descriptive encodings—both operator and
statistical—are ineffective or absent. The resulting behavior manifests
observationally as large-scale acceleration or repulsive effective dynamics.

\subsection{Relational encoding (E8)}

We recall the defining features of E8.

\begin{remark}
Relational encoding (E8) describes systems purely in terms of admissible
transitions and relations between configurations, without reference to local
state variables, fields, or ensemble distributions. It is the weakest descriptive
encoding in the MTT framework, but also the most robust.
\end{remark}

E8 persists when:
\begin{itemize}
\item local structure is absent or ill-defined;
\item statistical regularity is insufficient to define ensembles;
\item only global consistency constraints remain admissible.
\end{itemize}

\subsection{Loss of local and statistical encodings}

Dark energy regimes represent a further degradation of descriptive capacity.

\begin{remark}
In contrast to dark matter regimes, where statistical encoding (E9) remains
available, dark energy regimes are characterized by failure or decoupling of both
local encodings (E1--E6) and statistical encodings (E9). Only relational encoding
(E8) remains admissible.
\end{remark}

This represents the weakest nontrivial descriptive regime.

\subsection{Global consistency without local structure}

Relational encoding still constrains kinematic consistency.

\begin{lemma}
In a regime dominated by relational encoding, global transition structure
continues to influence kinematic consistency even in the absence of local
descriptive content.
\end{lemma}

\begin{proof}
Gravity encoding enforces kinematic consistency universally and acts on admissible
continuation relations themselves. Relational encoding supplies exactly such
relations, allowing global structure to affect admissible continuation without
local matter content.
\end{proof}

Thus dark energy influences large-scale kinematics despite lacking local degrees
of freedom.

\subsection{Effective repulsive behavior}

Relational dominance produces effective acceleration.

\begin{remark}
When only relational encoding remains, admissible continuation is constrained
primarily by global consistency rather than by local mass or energy content.
This can manifest as effective repulsive behavior or accelerated expansion in
geometric realizations.
\end{remark}

No negative-energy fields or exotic fluids are required.

\subsection{Homogeneity and isotropy}

Dark energy appears homogeneous.

\begin{remark}
Relational encoding lacks localized structure by definition. Its effects are
therefore naturally homogeneous and isotropic at large scales, consistent with
observational properties attributed to dark energy.
\end{remark}

This homogeneity is structural, not imposed.

\subsection{Persistence across scales}

Relational dominance is scale-sensitive but persistent.

\begin{remark}
At sufficiently large scales or late times, relational encoding may dominate as
local and statistical encodings become ineffective. Once dominant, E8 effects
persist uniformly because no finer local description is available to break them.
\end{remark}

This explains the cosmological nature of dark energy phenomena.

\subsection{Comparison with cosmological constant models}

We contrast this interpretation with standard models.

\begin{remark}
Cosmological constant models posit a new constant energy density. The relational-
encoding interpretation explains similar phenomenology without introducing new
constants, interpreting acceleration as a consequence of missing local
descriptions rather than additional energy content.
\end{remark}

This reframing shifts emphasis from dynamics to description.

\subsection{Structural advantages}

The relational-encoding interpretation has several advantages.

\begin{itemize}
\item It explains universality and homogeneity naturally.
\item It avoids fine-tuning of vacuum energy.
\item It introduces no new fields or particles.
\item It aligns with admissibility and kinematic consistency.
\end{itemize}

These features follow structurally.

\subsection{Interpretive caution}

We emphasize again the scope of the claim.

\begin{remark}
This interpretation does not assert that dark energy is literally relational
encoding. It asserts that dark energy phenomena are well explained by regimes in
which only relational encoding remains admissible. Other explanations are not
excluded, but are not required at the structural level.
\end{remark}

\subsection{Preview: observational consequences}

We now prepare to discuss observational implications.

\begin{remark}
In the next section we analyze observational consequences and constraints of the
dark sector interpreted as missing encodings, and contrast them with particle
and modified-gravity models.
\end{remark}

% ============================================================
% Next section will be:
% \section{Observational Consequences and Constraints}
% ============================================================
\section{Observational Consequences and Constraints}

In this section we discuss observational consequences of interpreting the dark
sector as arising from missing or ineffective encodings. The goal is not to
derive new cosmological models or predict numerical parameters, but to clarify
how known observational features are naturally accommodated and which classes of
observations are most sensitive to encoding availability.

\subsection{Gravitational coupling only}

A central observational feature of the dark sector is its apparent coupling
solely through gravity.

\begin{remark}
In the Modal Triplet Theory framework, gravity encoding enforces kinematic
consistency universally and remains operative in all descriptive regimes. By
contrast, gauge and local matter encodings may fail or decouple. This explains
why dark sector phenomena interact gravitationally while exhibiting no
electromagnetic, weak, or strong signatures.
\end{remark}

This behavior is structural and does not depend on new interaction postulates.

\subsection{Galaxy rotation curves and clustering}

Dark matter effects in galaxies and clusters are naturally reproduced.

\begin{remark}
In regimes dominated by statistical encoding (E9), aggregate distributions
contribute to gravitational consistency and may cluster under kinematic
constraints. Such clustering reproduces the effective mass distributions inferred
from galaxy rotation curves and gravitational lensing, without requiring
localized particle degrees of freedom.
\end{remark}

The success of standard dark matter phenomenology is therefore preserved.

\subsection{Large-scale structure formation}

Statistical encoding supports structure formation.

\begin{remark}
Statistical ensemble encodings admit large-scale correlations and hierarchical
clustering. These features are compatible with observed large-scale structure and
cosmic web formation, provided appropriate initial conditions and refinement
histories are realized.
\end{remark}

No modification of gravitational dynamics is required.

\subsection{Cosmic acceleration and homogeneity}

Dark energy observations are likewise accommodated.

\begin{remark}
Relational encoding (E8) lacks localized structure and therefore manifests
homogeneously at large scales. Its influence on kinematic consistency may appear
as effective accelerated expansion in geometric realizations, consistent with
supernova and cosmic microwave background observations.
\end{remark}

This behavior arises without introducing a cosmological constant or exotic
fields.

\subsection{Redshift dependence and cosmic epochs}

Encoding availability may vary with cosmic epoch.

\begin{remark}
The relative availability of local, statistical, and relational encodings may
change as refinement conditions evolve over cosmic time. This allows dark matter
and dark energy phenomena to dominate at different epochs without invoking
time-varying constants or new dynamical sectors.
\end{remark}

Such evolution is descriptive rather than dynamical.

\subsection{Consistency with precision cosmology}

The encoding interpretation is compatible with existing constraints.

\begin{remark}
Because the present framework does not alter gravitational encoding or introduce
new propagating degrees of freedom, it is compatible with precision tests of
General Relativity, cosmic microwave background anisotropies, and baryon acoustic
oscillation measurements at the level of effective behavior.
\end{remark}

Detailed numerical modeling belongs to realization-specific studies.

\subsection{Absence of direct detection signals}

The lack of non-gravitational detection is naturally explained.

\begin{remark}
Direct detection experiments presuppose local particle or field encodings. In
statistical- or relational-dominated regimes, such encodings are unavailable or
ineffective, rendering direct detection intrinsically unlikely rather than
merely experimentally challenging.
\end{remark}

This shifts interpretation from experimental failure to descriptive limitation.

\subsection{Distinction from modified gravity models}

We contrast with alternative approaches.

\begin{remark}
Modified gravity models alter the kinematic consistency encoding itself. The
present interpretation leaves gravity unchanged and attributes dark sector
phenomena to restricted encoding availability. This distinction leads to
different expectations in regimes where local encodings re-emerge.
\end{remark}

This provides a principled way to discriminate interpretations.

\subsection{Potential discriminators}

Although conservative, the framework suggests qualitative discriminators.

\begin{remark}
Observations sensitive to the emergence or breakdown of local descriptive
structure—such as environments with extreme refinement instability or horizon
effects—may distinguish between particle-based dark matter models and
encoding-regime interpretations.
\end{remark}

Such discriminators are structural rather than parametric.

\subsection{Interpretive restraint}

We reiterate the scope of observational claims.

\begin{remark}
The present analysis does not claim new quantitative predictions. It demonstrates
compatibility with existing observations and offers an alternative explanatory
framework. Detailed phenomenological modeling requires specific realization
choices and lies beyond the scope of this paper.
\end{remark}

\subsection{Preview: objections and alternatives}

We now prepare to address common objections.

\begin{remark}
In the next section we address potential objections to the encoding interpretation
of the dark sector and compare it with alternative explanatory frameworks.
\end{remark}

% ============================================================
% Next section will be:
% \section{Objections, Alternatives, and Clarifications}
% ============================================================
\section{Objections, Alternatives, and Clarifications}

In this section we address common objections to the interpretation of the dark
sector as arising from missing or ineffective encodings. We also clarify the
relationship between the present framework and alternative approaches, including
particle dark matter and modified gravity models.

\subsection{``Is this just modified gravity?''}

A frequent objection is that any explanation of dark sector phenomena without
new particles implicitly modifies gravity.

\begin{remark}
The present framework does not modify gravitational dynamics or the kinematic
consistency encoding identified with gravity. Gravity remains universal and
unchanged. Dark sector phenomena arise because gravity continues to operate even
when other encodings (local matter or gauge encodings) are unavailable or
ineffective.
\end{remark}

Thus, the explanation attributes dark sector effects to descriptive regimes, not
to altered gravitational laws.

\subsection{``Is this just an effective field theory in disguise?''}

Another objection is that statistical or relational encodings are merely coarse-
grained effective descriptions of underlying local degrees of freedom.

\begin{remark}
Effective field theories presuppose the existence of underlying local encodings
that are integrated out or averaged over. In contrast, the present framework
allows for regimes in which no admissible local encoding exists at all. In such
regimes, statistical or relational encodings are not approximations but the only
available descriptions.
\end{remark}

This distinction is structural rather than epistemic.

\subsection{``Does this deny the existence of dark matter particles?''}

The interpretation is sometimes misread as denying the existence of new particles.

\begin{remark}
The present analysis does not deny the existence of new particles or fields. It
asserts only that dark sector phenomena do not require them at the structural
level. Particle-based explanations remain possible but are not forced by
admissibility.
\end{remark}

This maintains compatibility with future discoveries.

\subsection{``Why does dark matter cluster like matter?''}

Dark matter exhibits clustering behavior similar to ordinary matter.

\begin{remark}
Statistical ensemble encoding (E9) admits stable aggregate distributions that
cluster under gravitational kinematic consistency. Clustering therefore reflects
gravitational organization of aggregate structure, not the presence of localized
particles.
\end{remark}

This explains clustering without invoking particle trajectories.

\subsection{``Why does dark energy appear uniform?''}

Dark energy appears homogeneous and isotropic.

\begin{remark}
Relational encoding (E8) lacks localized descriptive content by definition. Its
effects therefore manifest uniformly at large scales. Homogeneity is a
structural consequence of relational dominance, not an imposed symmetry.
\end{remark}

This matches observed large-scale uniformity.

\subsection{Comparison with particle dark matter models}

We contrast the present interpretation with particle-based models.

\begin{remark}
Particle dark matter models introduce new local degrees of freedom to explain
gravitational anomalies. The encoding-regime interpretation attributes the same
phenomena to absence of local descriptive encodings. Both approaches are
logically consistent; the latter is more economical at the structural level.
\end{remark}

The two approaches may be distinguished phenomenologically in regimes where local
encoding re-emerges.

\subsection{Comparison with modified gravity}

We also contrast with modified gravity theories.

\begin{remark}
Modified gravity models alter the kinematic consistency encoding itself. The
present framework leaves gravity intact and explains dark sector behavior through
restricted encoding availability. This leads to different expectations in strong-
field or high-resolution regimes.
\end{remark}

This distinction is principled and testable.

\subsection{``Is this unfalsifiable?''}

A common concern is falsifiability.

\begin{remark}
The encoding interpretation is falsifiable in the following sense: observation
of robust, non-gravitational interactions of dark sector phenomena would require
the existence of local descriptive encodings, contradicting the regime
assumptions. Conversely, persistent absence of such interactions supports the
interpretation.
\end{remark}

Falsifiability is therefore structural rather than parametric.

\subsection{Scope and restraint}

We restate the limits of the claims.

\begin{remark}
This paper does not propose a new cosmological model, does not fit data, and does
not predict numerical parameters. It provides an interpretive framework that
organizes existing observations in terms of encoding availability.
\end{remark}

This restraint is intentional.

\subsection{Preview: summary and outlook}

We now conclude the analysis.

\begin{remark}
In the final section we summarize the dark sector interpretation and place it
within the overall architecture of the Modal Triplet Theory Program.
\end{remark}

% ============================================================
% Next section will be:
% \section{Summary and Outlook}
% ============================================================
\section{Summary and Outlook}

In this paper we have reinterpreted dark matter and dark energy as phenomenological
manifestations of restricted encoding availability in the Modal Triplet Theory
framework. No new obstruction types, encoding responses, or dynamical principles
were introduced. Instead, dark sector phenomena were shown to arise naturally
when certain descriptive encodings fail or decouple, leaving only weaker but more
robust encodings active.

Dark matter was interpreted as a regime dominated by statistical ensemble
encoding (E9), in which aggregate structure remains admissible and
gravitationally active despite the absence of stable local descriptive
encodings. This interpretation explains the gravitational clustering of dark
matter, its apparent mass-like behavior, and the persistent absence of
non-gravitational interactions without postulating new particles or forces.

Dark energy was interpreted as a regime dominated by relational encoding (E8), in
which only global transition structure remains admissible. In such regimes,
gravitational kinematic consistency acts on relations rather than on local or
statistical matter content, producing effective large-scale acceleration,
homogeneity, and isotropy. This explanation avoids introducing a cosmological
constant or exotic energy component and instead attributes acceleration to
descriptive incompleteness.

Throughout, gravity played a special role as a universal encoding that remains
operative across all descriptive regimes. Gauge and local matter encodings were
shown to be regime-dependent and may fail or decouple without violating
structural consistency. This distinction explains why the dark sector couples
gravitationally while remaining invisible to non-gravitational probes.

The dark sector interpretation developed here is intentionally conservative. It
does not deny the possible existence of new particles, fields, or modifications
of gravity. Rather, it demonstrates that such additions are not structurally
required to explain observed dark sector phenomena. The interpretation is
falsifiable in a structural sense: robust detection of non-gravitational
interactions associated with dark sector effects would contradict the encoding
regime assumptions underlying this framework.

This paper completes the D-layer of the Modal Triplet Theory Program. Together
with the A-layer (structural core, coherent kinematics, predictive limits),
the B-layer (encoding responses and their coexistence), and the C-layer
(realizations), it demonstrates that a wide range of physical phenomena—gravity,
gauge structure, quantization, rigidity of the Standard Model, string-like
frameworks, and dark sector behavior—can be understood as consequences of the
limits of describability rather than as independent postulates.

Future work may pursue several directions. Detailed realization-dependent
cosmological models may be constructed to explore quantitative consequences of
encoding-regime transitions. Observational probes sensitive to the emergence or
breakdown of local descriptive structure may provide discriminators between
encoding-based and particle-based explanations. More broadly, the encoding
availability perspective may offer insight into other regimes where effective
descriptions fail or reorganize.

In this way, the Modal Triplet Theory Program provides a unified conceptual
framework in which both visible and dark aspects of the universe arise from the
same structural principles governing what can—and cannot—be described.

% ============================================================
% End of Paper D1

\end{document}
